\documentclass{article}
\usepackage{tikz}
\begin{document}
Inside the current text we place two pictures, containing nodes named \texttt{n1} and \texttt{n2}, using

\verb|\tikz[remember picture] \node[circle,fill=red!50] (n1) {};|

which yields \tikz[remember picture] \node[circle,fill=red!50] (n1) {}; , and

\verb|\tikz[remember picture] \node[fill=blue!50] (n2) {};|

yielding the code \tikz[remember picture] \node[fill=blue!50] (n2) {};. To connect these codes, we create another picture using the overlay option and also the \texttt{remember picture} option.

\begin{center}
\begin{verbatim}
\begin{tikzpicture}[remember picture,overlay]
\draw[->,very thick] (n1) -- (n2);
\end{tikzpicture}
\end{verbatim}
\end{center}

\begin{tikzpicture}[remember picture,overlay]
\draw[->,very thick] (n1) -- (n2);
\end{tikzpicture}

Note that the last picture is seemingly empty. What happens is that it has zero size and contains an arrow
that lies well outside its bounds. As a last example, we connect a node in another picture to the first two
nodes. Here, we provide the \texttt{overlay} option only with the line that we do not wish to count as part of the
picture.

\begin{tikzpicture}[overlay,remember picture]
\node (c) [circle,draw] {Big circle};
\draw [->,very thick,red,opacity=.5]
(c) to[bend left] (n1) (n1) -| (n2);
\end{tikzpicture}

\begin{tikzpicture}[remember picture]
\node (c) [circle,draw] {Big circle};
\draw [overlay,->,very thick,red,opacity=.5]
(c) to[bend left] (n1) (n1) -| (n2);
\end{tikzpicture}

\begin{figure}[h]
\centering
\begin{minipage}{0.4\linewidth}
\begin{tikzpicture}[remember picture]
\node (c) [circle,draw] {Big circle};
\draw [overlay,->,very thick,red,opacity=.5]
(c) to[bend left] (n1) (n1) -| (n2);
\end{tikzpicture}
\end{minipage}
\begin{minipage}{0.5\linewidth}
\begin{verbatim}
\begin{tikzpicture}[remember picture]
\node (c) [circle,draw] {Big circle};
\draw [overlay,->,very thick,red,opacity=.5]
(c) to[bend left] (n1) (n1) -| (n2);
\end{tikzpicture}
\end{verbatim}
\end{minipage}
\end{figure}

\clearpage
There is a special node called current page that can be used to access the \texttt{current page}. It is a node of shape rectangle whose \texttt{south west} anchor is the lower left corner of the page and whose north east anchor is the upper right corner of the page. While this node is handled in a special way internally, you can reference it as if it were defined in some remembered picture other than the current one. Thus, by giving the \texttt{remember picture} and the \texttt{overlay} options to a picture, you can position nodes \textit{absolutely} on a page.
The first example places some text in the lower left corner of the current page:

\begin{verbatim}
\begin{tikzpicture}[remember picture,overlay]
\node [xshift=1cm,yshift=1cm] at (current page.south west)
[text width=7cm,fill=red!20,rounded corners,above right]
{
This is an absolutely positioned text in the
lower left corner. No shipout-hackery is used.
};
\end{tikzpicture}
\end{verbatim}
\begin{tikzpicture}[remember picture,overlay]
\node [xshift=1cm,yshift=1cm] at (current page.south west)
[text width=7cm,fill=red!20,rounded corners,above right]
{
This is an absolutely positioned text in the
lower left corner. No shipout-hackery is used.
};
\end{tikzpicture}

The next example adds a circle in the middle of the page.
\begin{verbatim}
\begin{tikzpicture}[remember picture,overlay]
\draw [line width=1mm,opacity=.25]
(current page.center) circle (3cm);
\end{tikzpicture}
\end{verbatim}
\begin{tikzpicture}[remember picture,overlay]
\draw [line width=1mm,opacity=.25]
(current page.center) circle (3cm);
\end{tikzpicture}

The final example overlays some text over the page (depending on where this example is found on the page, the text may also be behind the page).
\begin{tikzpicture}[remember picture,overlay]
\node [rotate=60,scale=10,text opacity=0.2]
at (current page.center) {Example};
\end{tikzpicture}
\end{document}